\emph{(1)} The Majorana generators satisfy $e_j^2 = 1$ and anticommute. For $j$ outside the $k$-th plane, $B_k$ commutes with $e_j$ (two transpositions); for $j$ in it, $B_k e_{2k-1} B_k^{-1} = -e_{2k-1}$ and likewise for $e_{2k}$, with $B_k^{-1} = -B_k = \tilde B_k$. Since $B_k^2 = -1$, $R_k(\theta) = \exp(-\theta B_k/2) = \cos(\theta/2) - B_k\sin(\theta/2)$, and $R_k(\pi) = -B_k$, whose adjoint action is that of $B_k$: rotation by $\pi$ in the $k$-th plane. The product $B_1 \cdots B_n = e_1 e_2 \cdots e_{2n}$ is the volume element, and its adjoint action negates every generator, hence agrees with $\alpha$ on generators and so on the algebra. On Fock space, $J$ commuting with the mode operators, $-J e_{2k-1}e_{2k} = -J(a_k + a_k^*)\,J(a_k^* - a_k) = (a_k + a_k^*)(a_k^* - a_k) = a_k a_k^* - a_k^* a_k = 1 - 2N_k$, so $\prod_k(1 - 2N_k) = (-J)^n B_1\cdots B_n$ implements $\alpha$; this is the chirality of the spinor module of Fact \claimref{rungs.imported-fock-representation-bogoliubov-group}. In terms of the mode operators, $\operatorname{Ad}(e^{J\theta N_k})a_k = e^{-J\theta}a_k$, and $e^{J\pi N_k} = 1 - 2N_k$: the phase of a mode operator is the rotation angle of its plane.

\emph{(2)} $\alpha(H) = H$ means $\Gamma H \Gamma^{-1} = H$, so $\Gamma$ commutes with every polynomial in $H$ and, the space being finite dimensional, with every function of $H$ (a function of a diagonalisable matrix is a polynomial in it, the Lagrange interpolant on its spectrum). For $\rho = f(H)$, $\omega(\alpha(x)) = \operatorname{tr}(\rho\,\Gamma x \Gamma^{-1}) = \operatorname{tr}(\Gamma^{-1}\rho\Gamma\, x) = \operatorname{tr}(\rho x) = \omega(x)$.

\emph{(3)} A state $\rho$ is invariant under $\operatorname{Ad}(e^{JtH})$ for all real $t$ iff $[\rho, H] = 0$, iff $\rho$ is block diagonal with respect to the eigenspaces of $H$. Since $\Gamma$ commutes with $H$, each eigenspace $E_\lambda$ splits as $E_\lambda^+ \oplus E_\lambda^-$ into its even and odd parts. If every eigenspace has a definite parity (one part zero), every block-diagonal $\rho$ commutes with $\Gamma$, which is constant on each block, so it is even. If some $E_\lambda$ has both parts nonzero, take unit vectors $u^\pm \in E_\lambda^\pm$ and $\rho = \tfrac12(\lvert u^+\rangle + \lvert u^-\rangle)(\langle u^+\rvert + \langle u^-\rvert)$: it commutes with $H$, and $\operatorname{tr}(\rho\,\lvert u^+\rangle\langle u^-\rvert) = \tfrac12 \neq 0$ on an odd element, so it is not even. Both parts nonzero is exactly a common eigenvalue of $H$ on the two sectors. For $H = \sum_k\omega_kN_k$ the eigenvalues on the occupation basis are $\sum_{k\in T}\omega_k$ with parity $|T| \bmod 2$; a coincidence $\sum_{T}\omega = \sum_{T'}\omega$ with $|T|$, $|T'|$ of opposite parity is the vanishing of $\sum_{T\setminus T'}\omega_k - \sum_{T'\setminus T}\omega_k$, a signed sum over the symmetric difference, of odd size; conversely a vanishing signed sum over odd $S$ gives $T = \{\varepsilon_k = +1\}$, $T' = \{\varepsilon_k = -1\}$. If every $\omega_k = (2m_k+1)\omega_0$, a signed sum of an odd number of them is $\omega_0$ times a sum of an odd number of odd integers, which is odd, hence nonzero; and $e^{J\pi H/\omega_0} = \prod_k e^{J\pi(2m_k+1)N_k} = \prod_k(1-2N_k) = \Gamma$. When the spectra are disjoint, $\Gamma$ is the Lagrange polynomial in $H$ taking the value $\pm 1$ on the eigenvalues of the even and odd sectors, which the check confirms for weights $(1,3,5)$ and $(1,2)$.

\emph{(4)} For one mode of weight zero, $H = 0$ and every state is stationary; $\tfrac12(1 + \tfrac12\sigma_x)$ with $\sigma_x = a + a^*$ is positive and has $\omega(a + a^*) = \tfrac12$. For weights $(1,2,3)$, $x = a_1a_2a_3^*$ shifts the eigenvalue of $H$ by $-1-2+3 = 0$, so $X = x + x^*$ commutes with $H$; $X$ is Hermitian with $X^3 = X$ (the check confirms that $P = X^2$ is a projection with $XP = X$), so its eigenvalues lie in $\{-1,0,1\}$ and $\tfrac18(1 + \tfrac12 X)$ is a density matrix; it commutes with $H$ and not with $\Gamma$, and $\omega(X) = \tfrac{1}{16}\operatorname{tr}(X^2) \neq 0$ on the odd element $X$.

\emph{The final sentence.} The boost of the geometric modular condition (Definition \claimref{obj.singularity.geometric-modular-condition-g}) is $K_{\mathrm{geo}} = \sum_k \omega_k K_k$ with $K_k = \tfrac12\sigma_k = N_k - \tfrac12$ in the Fock representation and mode-dependent weights $\omega_k$. It is even, and it differs from $\sum_k\omega_kN_k$ by a constant, so its stationary states and its sector spectra are those of (3): every $K_{\mathrm{geo}}$-stationary state is even iff no signed sum of an odd number of the $\omega_k$ vanishes; a zero weight is such a sum on its own, and a common nonzero weight $\kappa$ is the case $\omega_0 = \kappa$ of (3). For (G): a faithful state $\rho$ on the finite-dimensional factor has modular group $\operatorname{Ad}(\rho^{Jt})$; if this equals $\operatorname{Ad}(e^{JtcK_{\mathrm{geo}}})$ for all $t$ (the constant $c$ absorbing the reparametrisation), then $\rho^{Jt}e^{-JtcK_{\mathrm{geo}}}$ commutes with the whole matrix algebra, so it is a scalar $e^{J t \gamma}$ (a continuous one-parameter group of scalars), and differentiating at $t = 0$ gives $\log\rho = cK_{\mathrm{geo}} + \gamma$: $\rho$ is a Gibbs state of $K_{\mathrm{geo}}$, even by (2). States outside these classes, and stationary states with a zero-weight mode or an odd resonance, are where the premise is an independent input, by (4).
